\documentclass[11pt,a4paper]{exam}
\usepackage[utf8]{inputenc}
\usepackage{amsmath, amsthm, amssymb, amsfonts} %ams packages
\usepackage{graphicx, color} %graphics
\usepackage{enumitem} %personalized lists 
\usepackage[top=3.5cm, bottom=2cm, left=2cm, right=2cm, headsep = 3pt, footskip=2pt]{geometry} %pagelayout
\usepackage{multirow,multicol} %to use multirow on tables
\usepackage{tabularx}
\usepackage{array} %bmatrix and similar stuff
% \usepackage{hyperref}
\usepackage[slovene]{babel}

\usepackage{multicol}


%%%%%%%%%%%%%%%%%%%%%%%%%%%%%%%%%%%%%%%%%%%%%%%%%%%%%%%%%%%%%%%%%%%%%
%personalized commands to make latex look better (optional)
%%%%%%%%%%%%%%%%%%%%%%%%%%%%%%%%%%%%%%%%%%%%%%%%%%%%%%%%%%%%%%%%%%%%%%%%%%%
\renewcommand{\leq}{\leqslant} %menores y mayores iguales decentes
\renewcommand{\geq}{\geqslant}
\renewcommand{\subset}{\subseteq} 
\renewcommand{\supset}{\supseteq} 
\renewcommand{\subsetneq}{\varsubsetneq}
\renewcommand{\{}{\lbrace}
\renewcommand{\}}{\rbrace}
\newcommand{\sm}{\setminus} %setminus corto

%\renewcommand{\bar}{\overline}
%\renewcommand{\hat}{\widehat}

%%%%%%%%%%%%%%%%%%%%%%%%%%%%%%%%%%%%%%%%%%%%%%%%%%%%%%%%%%%%%%%%%%%%%}

%%%%%%%%%%%%%%%%%%%%%%%%%%%%%%%%%%%%%%%%%%%%%%%%%%%%%%%%%%%%%%%%%%%%%%%%%%%
%exam class commands
%%%%%%%%%%%%%%%%%%%%%%%%%%%%%%%%%%%%%%%%%%%%%%%%%%%%%%%%%%%%%%%%%%%%%%%%%%%
\pagestyle{head}
% \chead{ Matematična analiza 2024/2025 \\ Univerza v Ljublani, Pedagoška fakulteta \\ Vaje 1 \\17. februar 2025} %self-explanatory
% \runningheader{}{ Abstraktna algebra 2024/2025 \\ Univerza v Ljublani, Pedagoška fakulteta \\2. Kolokvij \\20. januar 2025}{} %self-explanatory
% \runningheader{}{}{}
% \footer{}{}{}
% \firstpagefooter{MATH 1190 3.0 M  }{Page \thepage\ of \numpages}{Copyright 2020 \textcopyright Antonio Montero}
% \firstpagefooter{}{}{}
% \runningfooter{MATH 1190 3.0 M }{Page \thepage\ of \numpages}{Copyright 2020 \textcopyright Antonio Montero}
% \firstpagefootrule 
% \runningfootrule

%%Solutions
\unframedsolutions
% \printanswers
\SolutionEmphasis{\itshape}


%Personalized commands
\pointpoints{točka}{točk}
\bonuspointpoints{točka}{dodatnih točk}
% \gradetable[h][questions]
\addpoints
\hqword{Naloga}
% \vpgword{text}
\hpword{Točke}
% \hsword{}
% \hbpword{}
\htword{Skupaj}

% \pointsinrightmargin
% \extrawidth{-2cm}
\marginpointname{ \points}


\chqword{Naloga}
\chpgword{Stran}
\chpword{točke}
% \chbpword{Puntos extra}
% \chsword{Puntos obtenidos
\chtword{Skupaj}
%%%%%%%%%%%%%%%%%%%%%%%%%%%%%%%%%%%%%%%%%%%%%%%%%%%%%%%%%%%%%%%%%%%%%

% \graphicspath{{../img/}} %Uso: \graphicspath{{RelativePath}}

%%%%%%%%%%%%%%%%%%%%%%%%%%%%%%%%%%%%%%%%%%%%%%%%%%%%%%%%%%%%%%%%%%%%%
%Personalized commands

%Personalized commands

\usepackage{tabularx}
\newcolumntype{C}{>{\centering\arraybackslash}X}%
\newcolumntype{R}{>{\raggedleft\arraybackslash}X}%
\newcolumntype{L}{>{\raggedright\arraybackslash}X}%

\newcolumntype{\cc}[1]{>{\hsize=#1\hsize\raggedright\arraybackslash}X}%
\newcolumntype{\rr}[1]{>{\hsize=#1\hsize\raggedleft\arraybackslash}X}%
\newcolumntype{\ll}[1]{>{\hsize=#1\hsize\centering\arraybackslash}X}%

\newcommand{\TT}{\mathbf{T}}
\newcommand{\FF}{\mathbf{F}}

\DeclareMathOperator{\sh}{sh}
\DeclareMathOperator{\ch}{ch}
\let\th\relax
\DeclareMathOperator{\th}{th}
\DeclareMathOperator{\arsh}{arsh}
\DeclareMathOperator{\arch}{arch}
\DeclareMathOperator{\arth}{arth}
\DeclareMathOperator{\arctg}{arctg}

%%%%%%%%%%%%%%%%%%%%%%%%%%%%%%%%%%%%%%%%%%%%%%%%%%%%%%%%%%%%%%%%%%%%%
		\newif\ifmarking
% 		\markingtrue %comment to hide marking suggestions

		\newenvironment{marking}
		{\hrulefill
		 
		\textbf{Marking suggestions:} 
		\ttfamily }
		{\hrulefill }

\allowdisplaybreaks

% \renewcommand{\thepartno}{\roman{partno}}
\renewcommand{\partlabel}{(\textit{\thepartno})}
%%%%%%%%%%%%%%%%%%%%%%%%%%%%%%%%%%%%%%%%%%%%%%%%%%%%%%%%%%%%%%%%%%%%%

\begin{document}
\chead{Matematična analiza 2024/2025 \\ Univerza v Ljubljani, Pedagoška fakulteta \\ 2. Izpit \\1. julij 2025}
% \footer{}{}{}

% \vspace{1 cm}

\pagestyle{head}
	

	
\begin{questions}
\question 
	\begin{parts}	
		\part Pod kakšnim kotom se sekata grafa funkcij $f(x) = \sin(x)$ in $g(x)=\cos(2x)$?
		\part Pokažite, da se krivulji $x^{2} + y^{2} = 1$ in $xy=5$ sekata pravokotno.
	\end{parts}


  \question Iz kroga z polmerom $1$ izrežemo izsek in preostanek zvijemo v stožec. Kolikšen naj bo kot izseka, da bo imel dobljeni stožec največji prostornino?

  \question
  Izračunajte
  \begin{parts}
    \part obseg astroide, ki je podana z enačbo $x^{2/3} + y^{2/3} = 1$,
    \part obseg kardioide, ki je podana v polarnih koordinatah s predpisom $r(\theta)= 2(1+\cos \theta)$,
    \part dolžino grafa funkcije $f(x) = \frac{x^{2}}{4} - \frac{\ln x}{2}$ nad intervalom $[1, e]$.
  \end{parts}

  \question Dano je zaporedje funkcij $f_{n}: \mathbb{R} \to \mathbb{R}$, kjer je $f_{n}(x) = \frac{x}{1+n^{2}x^{2}}$.
  \begin{parts}
    
    \part Ugotovite, če  zaporedje $\{f_{n}\}$ konvergira po točkah in če konvergira, določite njegovo limitno funkcijo $f(x)$.
    \part Določite, če je zaporedje $\{f_{n}\}$ enakomerno konvergentno.
    \part Ali velja $f'(x) = \lim_{n \to \infty} f_{n}'(x)$ za vsak $x \in \mathbb{R}$?
  \end{parts}

  
    \question Naj bo $X$ poljubna neprazna množica. Označimo z $X^{\mathbb{N}}$ množico vseh zaporedij  $(x_{1}, x_{2}, \dots)$ točk iz $X$. Definirajmo razdaljo $d(x,y)$ dveh elementov $x=(x_{1}, x_{2}, \dots)$ in $y=(y_{1}, y_{2}, \dots)$ te množice takole: \[d(x,y) = \begin{cases} 0, & \text{če je } x=y;  \\ \frac{1}{n} & \text{če veljajo } x_{1}=y_{1}, x_{2}= y_{2}, \dots, x_{n-1} = y_{n-1}, \text{ in } x_{n} \neq y_{n}. \end{cases}\]
		\begin{parts}
			\part Dokažite, da res je $d$ metrika na množici $M$.
			\part Pokažite, da ima ta prostor tole lastnost: vsaka točka poljubne (odprte ali zaprte) krogle je središče te krogle; poleg tega sta vsaki dve krogli bodisi disjunktni ali pa ena od njiju vsebuje drugo.
  \end{parts}


\end{questions}

\end{document}
